\section{Introduction}

A central challenge for QCD, the gauge theory of the strong nuclear force, is the direct prediction of hadron structure. In particular, parton distribution functions (PDFs) and generalized parton distributions (GPDs) have, until recently, posed intractable difficulties for first principles' calculations. Although a complete, \textit{ab initio} determination of PDFs and GPDs has yet to appear, a promising new method was recently proposed in Ref.~\cite{Ji:2013dva}. In this approach, PDFs and GPDs are extracted from matrix elements of spatially extended operators between nucleon states at finite momentum. These matrix elements are generally referred to as quasi-PDFs or quasidistributions. Related frameworks have also been proposed; in Refs.~\cite{Ma:2014jla,Ma:2014jga,Ma:2017pxb} quasidistributions were treated as ``lattice cross sections'' from which light-front PDFs can be factorized, and Refs.~\cite{Radyushkin:2016hsy,Radyushkin:2017cyf,Orginos:2017kos} introduced and studied the closely related pseudodistributions.

Quasidistributions are matrix elements of time-local operators and therefore can be directly calculated using lattice QCD \cite{Carlson:2017gpk,Briceno:2017cpo}, in which QCD is discretized on a Euclidean hypercubic lattice. Preliminary nonperturbative results for both quasi- and pseudodistributions are encouraging \cite{Lin:2014zya,Alexandrou:2015rja,Chen:2016utp,Orginos:2017kos,Zhang:2017bzy,Alexandrou:2017qpu}. Moreover, a number of theoretical issues \cite{Li:2016amo,Rossi:2017muf,Carlson:2017gpk} have been clarified or solved, including a proof of the multiplicative renormalization of the spatial Wilson-line operator \cite{Ji:2017oey,Ishikawa:2017faj}, a proof of the factorization of light-front PDFs from quasidistributions \cite{Ma:2014jla,Ma:2014jga}, and a proof that the matrix element extracted from Euclidean correlation functions is identical to that defined through a Lehmann-Symanzik-Zimmermann reduction procedure in Minkowski space \cite{Briceno:2017cpo,Xiong:2017jtn,Ji:2017rah}.

One of the computational challenges that must be addressed is the presence of a power divergence induced in the quasi- and pseudodistributions by the finite lattice regulator. The Wilson-line operator, which defines these distributions, has a divergence that scales exponentially with the length of the Wilson line divided by the lattice spacing, and this divergence must be removed nonperturbatively. Several approaches have been suggested: the authors of Refs.~\cite{Chen:2016fxx} and \cite{Ishikawa:2016znu} proposed removing the power divergence through an exponentiated mass renormalization, and, more recently, regularization invariant momentum subtraction (RI/MOM \cite{Chen:2017mzz} and RI' \cite{Alexandrou:2017huk}) schemes have been suggested as nonperturbative renormalization procedures. In Ref.~\cite{Monahan:2016bvm}, we introduced an alternative approach, the smeared quasidistribution, which circumvents some of the challenges of the power divergence by taking advantage of the properties of the gradient flow \cite{Narayanan:2006rf,Luscher:2011bx,Luscher:2013cpa}.

The gradient flow is a classical evolution, or one-parameter mapping, of the original quark and gluon degrees of freedom in a new parameter, the flow time. The flow time exponentially suppresses UV field fluctuations, which corresponds to smearing out the original degrees of freedom in real space. The critical property of the gradient flow is that, up to a multiplicative fermion wave function renormalization \cite{Luscher:2013cpa}, finite correlation functions of the original theory remain finite at nonzero flow time \cite{Luscher:2011bx}. By fixing the flow time in physical units, one can ensure that matrix elements determined nonperturbatively at finite lattice spacing remain finite in the continuum limit. The gradient flow therefore provides a nonperturbative, gauge-invariant method to render the quasidistribution finite, even in the continuum \cite{Monahan:2016bvm}. The resulting continuum matrix element can then be related directly to the light-front PDF, or to the quasi or pseudodistributions renormalized in, for example, the $\ms$ scheme. In practice, it is simpler to carry out this matching procedure in several steps: first relate the smeared distribution to the distribution without smearing, and then use the known relation between the unsmeared distribution and the light-front PDF \cite{Xiong:2013bka,Ji:2015jwa,Wang:2017qyg,Stewart:2017tvs}. The perturbative calculation of the relevant matrix elements in the $\ms$ scheme appears in Ref.~\cite{Constantinou:2017sej}, and here my focus is the determination of the matrix elements at finite flow time.

I study the matrix element that defines the smeared quasidistribution at one loop in perturbation theory. Computing the matrix element of the smeared Wilson-line operator between external massless quark states at rest, enables all integrals to be evaluated analytically. At one loop, the Wilson-line power divergence manifests as a contribution linear in $\oz = z/r_\tau$, where $z$ is the length of the Wilson line and $r_\tau$ is the gradient flow smearing radius, for $\oz \gg 1$. Subtracting this contribution, the remaining matrix element is finite, and has a well-defined $\oz\to0$ limit. In the small flow-time regime, for which $\oz \gg 1$, the matrix element depends only logarithmically on $\oz$ and satisfies a relation analogous to the usual renormalization-group equation.

I start by briefly reviewing light-front PDFs, smeared quasidistributions and their relation in Sec.~\ref{sec:defns}. I discuss the perturbative calculation of the relevant matrix element in Sec.~\ref{sec:pertth}, provide some numerical results in Sec.~\ref{sec:numres}, and summarize in Sec.~\ref{sec:summary}.
